% Parte «theory» (clase section): la edita solo su dueño, declarado en roles.yaml.
% Escribe una oración por línea y separa los párrafos con una línea en blanco (02-convencion §4.3).
% El título de la sección y su etiqueta sec:theory están en tex/outline/theory.tex, del editor del artículo:
% aquí no van \section ni \chapter; usa \subsection y niveles menores (02-convencion REQ-CONV-94).

Esta parte presenta los conceptos, las definiciones y los modelos en los que se apoya el trabajo.
La notación común del documento está en \texttt{tex/notation.tex}.
